\section{Recipient exchanges at matched peak mass}
We permit one swap per query between an edited agent and a
previously unedited agent. The incoming path follows OLS4 and is
scaled down to the outgoing half-peak radius. Each query keeps
exactly the same number of edits and sum of per-row peak radii.
We compare a causal OLS4 geometry score with an exchange model
trained on analytic pairs, with or without requiring predicted
dual gains and non-increasing risk. Nine fits use
\ExchangeMassUpdates{} updates across three seeds and three
source-role rotations. The three exchange policies use the same
shortlist and the same \ExchangeMassEdits{} total edits as the
frozen half-peak anchor.

The geometry, learned and guarded policies make
\ExchangeMassGeometryCount{}, \ExchangeMassLearnedCount{} and
\ExchangeMassGuardCount{} exchanges, respectively.
Their paired conservative utility contrasts against no exchange
are \ExchangeMassGeometryEffect{}, \ExchangeMassLearnedEffect{}
and \ExchangeMassGuardEffect{} pp. Their paired distance
contrasts are \ExchangeMassGeometryPair{}, \ExchangeMassLearnedPair{}
and \ExchangeMassGuardPair{} pixels. Intervals use paired source
resampling. Original recording risk passes are
\ExchangeMassPasses{} out of 72 (anchor, geometry, learned, guard).
The worst signed easy changes in that order are
\ExchangeMassWorst\%. Calibration qualifies \ExchangeMassQual{}/9.

Ignoring risk and maximizing each conservative gain separately
over the same exchange menu excludes OLS4 gains in
\ExchangeMassOLSLimit{} of 72 views. This bound does not
establish that the remaining views admit a joint positive
and risk-feasible choice. Unlike unrestricted recipient
reallocation, it preserves the anchor query counts and peak
mass and permits only one exchange in each query.
No independent confirmation outcomes enter this comparison.

We also relax the menu completely. Let $d_{it}$ be CV error on
known steps and CV-to-OLS4 disagreement on missing steps. Giving
row $i$ radius $r$ can improve the conservative OLS4 numerator
by at most $g_i(r)=w_i\sum_t\min(r,d_{it})$. For total radius $B$
and at most $K$ edits, any $\lambda\geq0$ gives the upper bound
\[
 \lambda B + \operatorname{TopK}_{i}\sup_{r\geq0}
          \bigl(g_i(r)-\lambda r\bigr).
\]
Here $\operatorname{TopK}$ sums the $K$ largest row values.
We minimize over a price grid and add the original CV-to-OLS4
offset. Query-local budgets bind each query separately; recording
pooling keeps only the total radius and count per recording.
Both allow arbitrary pointwise paths and ignore risk.

\begin{center}
\small
\setlength{\tabcolsep}{3pt}
\begin{tabular}{@{}lrrr@{}}
\toprule
Peak factor & Query fixed & Query free & Record fixed \\
\midrule
$0.5\times$ & \ExchangeMassHalfLimits \\
$1\times$ & \ExchangeMassBaseLimits \\
$2\times$ & \ExchangeMassDoubleLimits \\
$4\times$ & \ExchangeMassFourLimits \\
$8\times$ & \ExchangeMassEightLimits \\
Unbounded & \ExchangeMassFreeLimits \\
\bottomrule
\end{tabular}
\end{center}
Each entry is the number of OLS4 gains ruled out of 72,
relative to the half-peak anchor budget. Free query count removes
the count restriction, not the radius constraint. Even after
pooling the same total radius and count across queries, the bound
rules out \ExchangeMassRecordExcluded{} views. This distinguishes
the aggregate displacement bottleneck from candidate ordering.
